


\chapter{Exercicis proposats}

\begin{enumerate}
\item Demostreu que el conjunt $\mathbb{C}$ dels nombres complexos,
amb les operacions%
\begin{align*}
(a+bi)+(c+di) &=(a+c)+(b+d)i \\
\lambda (a+bi) &=\lambda a+\lambda bi
\end{align*}%
té estructura d’espai vectorial real.

\item Considereu el conjunt de les funcions definides sobre l'interval $%
[0,1]$ i tals que $2f(0)=f(1)$. Formen un espai
vectorial real?

\item Considereu el conjunt dels polinomis de grau menor o igual a
4, sobre un cos $\mathbb{K}$ (inclòs el polinomi nul).
Formen un espai vectorial sobre $\mathbb{K}$?

\item Considerem l’espai vectorial real $\mathcal{F}(\mathbb{R},\mathbb{%
R})$ de les funcions reals de variable real. Determineu quins dels
següents subconjunts són subespais vectorials: (a) $S_{1}=\{f\in 
\mathcal{F}(\mathbb{R},\mathbb{R}):f(-x)=f(x),\forall x\in \mathbb{R}\}$;
(b) $S_{2}=\{f\in \mathcal{F}(\mathbb{R},\mathbb{R}):f(-x)=-f(x),\forall
x\in \mathbb{R}\}$; (c) $S_{3}=\{f\in \mathcal{F}(\mathbb{R},\mathbb{R}):$ $%
f $ és contínua$\}$; (d) $S_{4}=\{f\in \mathcal{F}(\mathbb{R},\mathbb{R}%
):f(x)\geq 0,\forall x\in \mathbb{R}\}$; (e) $S_{5}=\{f\in \mathcal{F}(%
\mathbb{R},\mathbb{R}):f$ és dues vegades derivable i $f^{\prime \prime
}-f^{\prime }+f=\mathbf{0}\}$; (f) $S_{6}=\{f\in \mathcal{F}(\mathbb{R},%
\mathbb{R}):f(1)=f(2)+3\}$.

\item Considerem l’espai vectorial real $\mathcal{M}_{n\times m}(%
\mathbb{R})$ de les matrius d’ordre $n\times m$ amb coeficients reals.
Determineu quins dels següents subconjunts són subespais
vectorials: (a) $S_{1}=\{A\in \mathcal{M}_{n\times m}(\mathbb{R}%
):a_{11}=0\} $; (b) $S_{2}=\{A\in \mathcal{M}_{n\times m}(\mathbb{R}%
):\sum_{i=1}^{n}a_{ii}=0,n=m\}$; (c) $S_{3}=\{A\in \mathcal{M}_{n\times m}(%
\mathbb{R}):a_{ij}=a_{ji},\forall i,j,n=m\}$.

\item Considérese les inclusiones $\mathbb{Q}\subset \mathbb{Q}(\sqrt[3]{%
5})\subset \mathbb{R}\subset \mathbb{C}$, essent $\mathbb{Q}(\sqrt[3]{5}%
)=\{a+b\sqrt[3]{5}:a,b\in \mathbb{Q}\}$. Demostreu que cadascun d’aquests
cos és un espai vectorial sobre l'anterior.

\item Escriviu el vector $\mathbf{w}=(2,2,3)$ com una combinació
lineal dels vectors del sistema $\mathcal{S}=(\mathbf{v}_{1},\mathbf{v}%
_{2},\mathbf{v}_{3})$ de $\mathbb{R}^{3}$, essent $\mathbf{v}_{1}=(2,1,4),%
\mathbf{v}_{2}=(1,-1,3),\mathbf{v}_{3}=(3,2,5)$.

\item Escriviu el polinomi $p(x)=2+2x+3x^{2}$ com una combinació
lineal dels polinomis $%
p_{1}(x)=2+x+4x^{2},p_{2}(x)=1-x+3x^{2},p_{3}(x)=3+2x+5x^{2}$ de $\mathbb{R}%
_{2}\left[ x\right] $.

\item Esbrineu si la matriu%
\[
A=\left( 
\begin{array}{cc}
-1 & 7 \\ 
5 & 1%
\end{array}%
\right) 
\]%
és o no una combinació lineal de les matrius%
\[
M_{1}=\left( 
\begin{array}{cc}
1 & 2 \\ 
-1 & 3%
\end{array}%
\right) ,M_{2}=\left( 
\begin{array}{cc}
0 & 1 \\ 
2 & 4%
\end{array}%
\right) ,M_{3}=\left( 
\begin{array}{cc}
4 & -2 \\ 
0 & -2%
\end{array}%
\right) 
\]%
de $\mathcal{M}_{2}(\mathbb{R})$.

\item (a) Expresseu $(4a,a-b,a+2b)$ com una combinació lineal de $%
(4,1,1) $ i $(0,-1,2)$. (b) Expresseu $(3a+b+3c,-a+4b-c,2a+b+2c)$ com una
combinació lineal de $(3,-1,2)$ i $(1,4,1)$. (c) Expresseu $%
(2a-b+4c,3a-c,4b+c)$ com una combinació lineal de tres vectors
diferents de zero.

\item Expresseu $\mathbf{v}=(1,1)$ com una combinació lineal de $\mathbf{%
v}_{1}=(1,-1),\mathbf{v}_{2}=(3,0),\mathbf{v}_{3}=(2,1)$, de dues maneres
diferents.

\item Expresseu (si és possible) els següents elements de $\mathcal{F}(%
\mathbb{R},\mathbb{R})$ com una combinació lineal de les famílies
corresponents: (a) $\sin x$ i $(1,x,x^{2},...)$; (b) $x^{2}+x-1$ i $%
(1,x-1,(x-1)^{2})$; (c) $1$ i $(x+1,x^{2}-1,x^{3}+1)$; $0$ i $%
((x-1)^{2},x,x^{2}+2,\sqrt{3})$.

\item En l’espai vectorial $\mathbb{R}^{4}$ es considera el subespai
generat pels vectors $(2,3,1,-5)$ i $(0,2,-1,3)$. Trobeu $a$ i $b$
perquè el vector $(2,a,3,-b)$ pertanyi a aquest subespai.

\item Trobeu un sistema generador $\mathcal{S}$ de l’espai vectorial de les
solucions dels següents sistemes lineals homogenis:%
\[
\text{(a) \ }\left. 
\begin{array}{r}
x+4y+8z=0 \\ 
2x+5y+6z=0 \\ 
3x+i-4z=0%
\end{array}%
\right\} \text{ \ (b) \ }\left. 
\begin{array}{r}
2x-3y+z=0 \\ 
6x-9y+3z=0 \\ 
-4x+6y-2z=0%
\end{array}%
\right\} 
\]

\item Dels sistemes següents, determineu quins engendran $\mathbb{R}%
^{3}$. En cas negatiu, doneu una descripció geomètrica del
subespai que engendran. (a) $\mathcal{S}_{1}=((4,7,3),(-1,2,6),(2,-3,5))$;
(b) $\mathcal{S}_{2}=((-2,5,0),(4,6,3))$; $\mathcal{S}%
_{3}=((1,0,3),(2,0,-1),(4,0,5),(2,0,6))$. 

\item Trobeu un sistema d'equacions lineals homogènies el qual espai
de solucions és cadascun dels següents subespais vectorials:%
\begin{align*}
S_{1} &=\left\langle
(1,-2,0,3),(1,-1,-1,4),(1,0,-2,5)\right\rangle  \\
S_{2} &=\left\langle (1,2,0,1)\right\rangle  \\
S_{3} &=\left\langle 8+3x,2x+3x^{2}\right\rangle \text{ \
(considerat com a subespai de }\mathbb{R}_{2}\left[ x\right] \text{)} \\
S_{4} &=\left\langle \left( 
\begin{array}{cc}
0 & 1 \\ 
3 & 1%
\end{array}%
\right) ,\left( 
\begin{array}{cc}
0 & 1 \\ 
1 & 3%
\end{array}%
\right) ,\left( 
\begin{array}{cc}
3 & 1 \\ 
1 & 0%
\end{array}%
\right) \right\rangle 
\end{align*}

\item Es consideren els subespais $S$ i $T$ de $\mathbb{R%
}^{4}$ definits de la manera següent:%
\begin{align*}
S &=\{(x,y,z,t)\in \mathbb{R}^{4}:x=y=z\} \\
T &=\{(\alpha +\beta +\gamma +3\delta ,\alpha +\beta +2\delta
,\alpha +\delta ,\alpha +\gamma ):\alpha ,\beta ,\gamma ,\delta \in \mathbb{R%
}\}
\end{align*}%
(a)\ Està $T\subset S$? (b) Calculeu 
$S\cap T$.

\item Trobeu $S\cap T$ i $S+T$, essent%
\begin{align*}
S &=\left\langle (1,0,1),(1,1,0)\right\rangle  \\
T &=\left\langle (1,2,3),(0,0,1)\right\rangle 
\end{align*}%
Esbrineu si la suma és o no directa. Esbrineu també si $S$
i $T$ són o no subespais suplementaris respecto $\mathbb{R}^{3}$.

\item Demostreu que els sistemes $\mathcal{S}_{1}=((1,2,-1),(2,-1,0))$ i $%
\mathcal{S}_{2}=((-4,7,-2),(-1,-7,3))$ de $\mathbb{R}^{3}$ són equivalents.

\item En l’espai vectorial $\mathbb{R}^{4}$ es consideren els vectors $%
(1,2,0,a),(3,2,b,-1),(2,7,a,-3)$. Determineu $a$ i $b$ perquè aquests
vectors siguin linealment dependents.

\item En $\mathbb{R}^{4}$ estudiar la independència lineal dels vectors 
\[
(1,7,-1,0),(2,13,0,1),(3,17,4,2),(0,5,-2,1)
\]

\item Proveu que si $u_{1},u_{2},u_{3}$ són vectors linealment
independents d’un espai vectorial $E$, també ho són $%
u_{1}+u_{2},u_{2}+u_{3},u_{1}+u_{3}$.

\item Calculeu l’espai vectorial quocient de l’espai vectorial $\mathbb{R%
}^{4}$ mòdul el subespai $F$, essent: (a) $F%
=\left\{ (x,y,z,t):x+y=z-t=0\right\} $; (b) $F=\left\langle
(1,1,1,1),(1,2,3,4),(1,3,5,6)\right\rangle $; (c) $F%
=\{(x,y,z,t):x+y=x-y=z+t=x-t=0\}$.

\item Siguin 
\[
M=\{\left( 
\begin{array}{cc}
a & b \\ 
0 & c%
\end{array}%
\right) :a,b,c\in \mathbb{R}\}
\]%
\[
M^{\prime }=\{\left( 
\begin{array}{cc}
a-b & 2a \\ 
b & a+2b%
\end{array}%
\right) :a,b\in \mathbb{R}\}
\]%
Trobeu una base de $\mathcal{M}_{2\times 2}(\mathbb{R})/M$ i $\mathcal{M}%
_{2\times 2}(\mathbb{R})/M^{\prime }$.

\item Determineu els rangs dels sistemes de vectors de $\mathbb{R}^{4}$ i, si escau,
les relacions entre els vectors de $\mathcal{S}_{i}$ ($i=1,2$). (a) $%
\mathcal{S}_{1}:\mathbf{v}_{1}=(1,0,1,0),\mathbf{v}_{2}=(2,1,0,1),\mathbf{v}%
_{3}=(0,2,-1,1),\mathbf{v}_{4}=(3,-1,2,0)$; (b) $\mathcal{S}_{2}:\mathbf{v}%
_{1}=(1,0,2,3),\mathbf{v}_{2}=(7,4,-2,-1),\mathbf{v}_{3}=(5,2,4,7),\mathbf{v}%
_{4}=(3,2,0,1)$.

\item Determineu, segons els valors de $\alpha $ (o de $\alpha ,\beta
,\gamma $) els rangs dels sistemes $\mathcal{S}_{1},\mathcal{S}_{2},%
\mathcal{S}_{3}$ de vectors de $\mathbb{R}^{3}$. (a) $\mathcal{S}_{1}:%
\mathbf{v}_{1}=(\alpha ,1,1),\mathbf{v}_{2}=(1,\alpha ,1),\mathbf{v}%
_{3}=(1,1,\alpha )$; (b) $\mathcal{S}_{2}:\mathbf{v}_{1}=(\alpha ,1,1),%
\mathbf{v}_{2}=(-1,-\alpha ,-1),\mathbf{v}_{3}=(-1,-1,\alpha )$; $\mathcal{S}%
_{3}:\mathbf{v}_{1}=(0,\gamma ,-\beta ),\mathbf{v}_{2}=(-\gamma ,0,\alpha ),%
\mathbf{v}_{3}=(\beta ,-\alpha ,0)$.

\item Determineu una base del subespai $F$ de $\mathbb{R}^{4}$
descrit per $(x,y,z,t)$ del qual sap que $x=y-3z$ i $z=t$. Completeu la
base obtinguda per a obtenir una base de $\mathbb{R}^{4}$.

\item En $\mathbb{R}^{3}$ verificar que els vectors $\mathbf{a},\mathbf{b},%
\mathbf{c}$ són independents i calculeu les coordenades de $\mathbf{x}$
respecte de la base $(\mathbf{a},\mathbf{b},\mathbf{c})$. Els vectors són $%
\mathbf{a}=(-1,1,1),\mathbf{b}=(1,-1,1),\mathbf{c}=(1,1,-1)$ i $\mathbf{x}%
=(2,-3,-1)$.

\item Determineu el rang del sistema de vectors $\mathbf{v}_{i}$ ($%
i=1,...,n$) de $\mathbb{R}^{n}$ que tenen coordenades $\alpha
_{ij}=1+(\alpha -1)\delta _{ij}$, on $\alpha $ és un escalar conegut
i $\delta _{ij}$ és el símbol de Kronecker.

\item Siguin $F,G,H$ subespais de l’espai vectorial $E$. Demostreu o doneu
contraexemples de les afirmacions següents:

a) $F\cap (G+H)=(F\cap G)+(F\cap H)$

b) $F+(G\cap H)=(F+G)\cap (F+H)$

c) $\dim (F\cap (G+H))=\dim (F\cap G)+\dim (F\cap H)+\dim (F\cap G\cap H)$

\item Considerem l’espai vectorial $\mathbb{R}_{3}\left[ t\right] $ de
els polinomis amb coeficients reals de grau com a màxim tres. Siguin $F$ i $G
$ els subespais

$F=\left\{ p(t)\in \mathbb{R}_{3}\left[ t\right] :p(0)=p(1)=p^{\prime
}(1/2)=p^{\prime \prime \prime }(0)=0\right\} $

$G=\left\langle t+1,t,t-1\right\rangle $

a) Calculeu les seves dimensions i trobeu una base de cadascun d'ells

b) Considerem els polinomis $t^{2}-5t+2$ i $3t-4$. Es demana: (i) 
Són de $G$? (ii) En cas afirmatiu, trobeu $\alpha ,\beta
,\gamma \in \mathbb{R}$ tals que%
\[
p(t)=\alpha (t+1)+\beta t+\gamma (t-1)
\]%
essent $p(t)$ cadascun dels dos polinomis anteriors. (iii) 
Són aquests $\alpha ,\beta ,\gamma $ únics? En cas
afirmatiu, justificar la resposta. En cas contrari, trobar-los.

c) Calculeu les dimensions i trobeu una base de $F+G$ i $F\cap G$. 
És la suma $F+G$ directa?

d) Trobeu una base d’un complementari de $F+G$ en $\mathbb{R}_{3}\left[ t%
\right] $.

\item Considerem en $\mathbb{C}^{3}$ el sistema $\mathcal{S}=\left(
(2-i,1,1),(1,2+i,3),(1,2+3i,2)\right) $. Esbrineu: (a) Si $\mathcal{S}$ és
un sistema generador d’un subespai de $\mathbb{C}^{3}$ de dimensió
1; (b) Si $\mathcal{S}$ és un sistema generador d’un subespai de $\mathbb{%
C}^{3}$ de dimensió 2; (c) És una base de $\mathbb{C}%
^{3}$? (d) És un sistema lliure?
\end{enumerate}

\section{Pràctica amb ordinador}

Els càlculs d'aquesta unitat (rangs, independència, bases, interseccions) es
poden comprovar amb \textbf{SageMath}, un sistema de càlcul simbòlic lliure
que treballa amb aritmètica exacta sobre $\mathbb{Q}$ i sobre cossos finits.
No cal instal·lar res: els fragments següents es poden executar al navegador
a \texttt{https://sagecell.sagemath.org}. Us proposem de resoldre primer cada
exercici a mà i fer servir l'ordinador només per verificar el resultat.

\begin{enumerate}
\item Comproveu el rang del sistema de vectors de l'exercici corresponent
del capítol de pràctica:
\begin{verbatim}
A = matrix(QQ, [[1,1,2,0], [1,1,0,6], [-1,2,0,1], [1,1,1,3]])
print(A.rank())          # rang del sistema
print(A.kernel())        # relacions de dependència entre files
\end{verbatim}

\item Calculeu una base de la intersecció de dos subespais de
$\mathbb{R}^{4}$ i compareu-la amb el càlcul manual:
\begin{verbatim}
V = QQ^4
S = V.span([[1,2,1,0], [-1,1,1,1]])
T = V.span([[2,-1,0,1], [1,-1,3,7]])
I = S.intersection(T)
print(I.dimension()); print(I.basis())
\end{verbatim}

\item Sobre el cos finit $\mathbb{Z}_{5}$, estudieu la dependència lineal
del sistema de l'exercici de sistemes homogenis:
\begin{verbatim}
V = GF(5)^3
S = V.span([[1,4,0], [2,1,1], [1,1,4]])
print(S.dimension())     # dimensió del subespai generat
\end{verbatim}
Observeu que un mateix sistema de coeficients pot tenir rang diferent
segons el cos sobre el qual es treballa.
\end{enumerate}